\documentclass{article}
\usepackage{amssymb}
\usepackage{amsmath}
\usepackage{amsthm}
\usepackage{tikz-cd}
\usepackage{mathpartir}
\usepackage{stmaryrd}
\usepackage{parskip}
\usepackage{xcolor}
\usepackage{array}
\usepackage[margin=1.0in]{geometry}
\usepackage{comment}
\usepackage{quiver}

\definecolor{myorange}{RGB}{242, 114, 0}
\definecolor{myblue}{RGB}{0,162,255}



%%%%%%%%%%%%%%%%%%%%%%%%
% theorem environments
%%%%%%%%%%%%%%%%%%%%%%%%
\newtheorem{theorem}{Theorem}[section]
\newtheorem{lemma}{Lemma}[section]
\newtheorem{nonnum-theorem}{Theorem}[section]
\newtheorem{corollary}{Corollary}[section]

\theoremstyle{definition}
\newtheorem{definition}{Definition}[section]

\theoremstyle{remark}
\newtheorem{remark}{Remark}

\newtheorem{example}{Example}
%%%%%%%%%%%%%%%%%%%%%%%%%%%%%%%%%%%%%%%%%%%%%%%%%%%

\newcommand{\To}{\Rightarrow}
\newcommand{\inl}{\mathsf{inl}}
\newcommand{\inr}{\mathsf{inr}}
\newcommand{\alt}{\mathrel{\bf \,\mid\,}}

\newcommand{\ob}{\text{Ob}}


\newcommand{\extlc}{\text{Ext-}\lambda}
\newcommand{\extlcm}{\text{Ext-}\lambda^{-\text{trans}}}
\newcommand{\extlcmm}{\text{Ext-}\lambda^{-\text{trans}-\text{cast}}}
\newcommand{\extlcprime}{\text{Ext-}\lambda'}
\newcommand{\intlc}{\text{Int-}\lambda}
\newcommand{\intlcbisim}{\text{Int$_\approx$-}\lambda}
\newcommand{\erase}[1]{\lfloor {#1} \rfloor}


\newcommand{\uarrowl}{\mathrel{\rotatebox[origin=c]{-30}{$\leftarrowtail$}}}
\newcommand{\uarrowr}{\mathrel{\rotatebox[origin=c]{60}{$\leftarrowtail$}}}
\newcommand{\darrowl}{\mathrel{\rotatebox[origin=c]{30}{$\twoheadleftarrow$}}}
\newcommand{\darrowr}{\mathrel{\rotatebox[origin=c]{120}{$\twoheadleftarrow$}}}
\newcommand{\vuarrow}{\mathrel{\rotatebox[origin=c]{-90}{$\leftarrowtail$}}}
\newcommand{\vdarrow}{\mathrel{\rotatebox[origin=c]{90}{$\twoheadleftarrow$}}}

% Types, terms, and precision 

\newcommand{\dyn}{?}
\newcommand{\nat}{\text{Nat}}
\newcommand{\bool}{\text{Bool}}
\newcommand{\ra}{\rightharpoonup}
\newcommand{\Ret}[1]{\mathsf{Ret\,}{#1}}
\newcommand{\hole}[1]{\bullet \colon {#1}}
\newcommand{\dyntodyn}{\dyn \ra\, \dyn}


\newcommand{\up}[2]{\langle{#2}\uarrowl{#1}\rangle}
\newcommand{\dn}[2]{\langle{#1}\darrowl{#2}\rangle}

\newcommand{\upc}[1]{\text{up}\,{#1}\,}
\newcommand{\dnc}[1]{\text{dn}\,{#1}\,}


% \newcommand{\ret}{\mathsf{ret}}
\newcommand{\err}{\mho}
\newcommand{\zro}{\textsf{zro}}
\newcommand{\suc}{\textsf{suc}}
\newcommand{\lda}[2]{\lambda {#1} . {#2}}

\newcommand{\injarr}[1]{\textsf{Inj}_\ra ({#1})}
\newcommand{\injnat}[1]{\textsf{Inj}_\text{nat} ({#1})}
\newcommand{\casenat}[4]{\text{Case}_\text{nat} ({#1}) \{ \text{no} \to {#2} \alt \text{nat}({#3}) \to {#4} \}}
\newcommand{\casearr}[4]{\text{Case}_\ra ({#1}) \{ \text{no} \to {#2} \alt \text{fun}({#3}) \to {#4} \}}
\newcommand{\casedyn}[5]{\text{Case}_\Dyn ({#1}) \{ \text{nat}({#2}) \to {#3} \alt \text{fun}({#4}) \to {#5} \}}
%\newcommand{\bind}[3]{\text{var } {#1} = {#2} \text{ in } {#3}}
\newcommand{\matchnat}[4]{\text{match } {#1} \text{ with } \{ \textsf {zero} \Rightarrow {#2} \alt \textsf{ suc } {#3} \Rightarrow {#4} \}}

\newcommand{\refl}{\text{refl}}

\newcommand{\Lift}{\text{Lift}}

%\newcommand{\rel}{\circ\hspace{-4px}-\hspace{-4px}\bullet}
\newcommand{\rel}{\mathrel{\circ\mkern-6mu-\mkern-6mu\bullet}}
% \newcommand{\wand}{\mathrel{-\mkern-6mu*}}
\newcommand{\wand}{\mathrel{\multimap}}
\newcommand{\ltdyn}{\sqsubseteq}
\newcommand{\gtdyn}{\sqsupseteq}
\newcommand{\equidyn}{\mathrel{\gtdyn\ltdyn}}
\newcommand{\gamlt}{\Gamma^\ltdyn}
\newcommand{\deltalt}{\Delta^\ltdyn}
\newcommand{\relcomp}{\odot}

\newcommand{\hasty}[3]{{#1} \vdash {#2} \colon {#3}}
\newcommand{\vhasty}[3]{{#1} \vdash^v {#2} \colon {#3}}
\newcommand{\phasty}[3]{{#1} \vdash^c {#2} \colon {#3}}
\newcommand{\etmprec}[4]{{#1} \vdash {#2} \ltdyn_e {#3} \colon {#4}}
\newcommand{\itmprec}[4]{{#1} \vdash {#2} \ltdyn_i {#3} \colon {#4}}
\newcommand{\etmequidyn}[4]{{#1} \vdash {#2} \equidyn_e {#3} \colon {#4}}
\newcommand{\itmequidyn}[4]{{#1} \vdash {#2} \equidyn_i {#3} \colon {#4}}
\newcommand{\synbisim}{\approx_{\text{syn}}}

\newcommand{\Dwn}{\Downarrow}
\newcommand{\qte}[1]{\text{quote}({#1})}

\newcommand{\elab}[1]{\text{Elab}({#1})}

% Perturbations
\newcommand{\pertp}{\text{Pert}^\text{P}}
\newcommand{\perte}{\text{Pert}^\text{E}}
\newcommand{\pertdyn}[2]{\text{pert-dyn}({#1}, {#2})}
\newcommand{\delaypert}[1]{\text{delay-pert}({#1})}

\newcommand{\pertc}{\text{Pert}_{\text{C}}}
\newcommand{\pertv}{\text{Pert}_{\text{V}}}


% SGDT and Intensional Stuff

\newcommand{\later}{{\vartriangleright}}
\newcommand{\laterhs}{{\later}}
\newcommand{\type}{\texttt{Type}}
\newcommand{\lob}{\text{L\"{o}b}}
\newcommand{\tick}{\mathsf{tick}}
\newcommand{\nxt}{\mathsf{next}}
\newcommand{\fix}{\mathsf{fix}}
\newcommand{\kpa}{\kappa}

% Model-related stuff
\newcommand{\calV}{\mathcal{V}}
\newcommand{\calE}{\mathcal{E}}
\newcommand{\calVk}{\mathcal{V}_k}
\newcommand{\calEk}{\mathcal{E}_k}
\newcommand{\tok}{\mathrel{\mathop{\to}\limits^{\textrm{\footnotesize k}}}}
\newcommand{\timesk}{\mathrel{\mathop{\times}\limits^{\textrm{k}}}}
\newcommand{\Fk}{\mathrel{\mathop{F}\limits^{\textrm{k}}}}
\newcommand{\Uk}{\mathrel{\mathop{U}\limits^{\textrm{k}}}}

\newcommand{\op}[1]{{#1}^{\textrm{op}}}
\newcommand{\calC}{\mathcal{C}}
\newcommand{\Set}{\mathsf{Set}}
\newcommand{\ErrDom}{\mathsf{ErrDom}}
\newcommand{\Yo}{\mathsf{Yo}}
\newcommand{\Hom}{\mathsf{Hom}}
\newcommand{\calS}{\mathcal{S}}
\newcommand{\gfix}{\texttt{gfix}}
\newcommand{\qfix}{\texttt{qfix}}
\newcommand{\calU}{\mathcal{U}}
\newcommand{\laterhat}{\widehat{\later}}
\newcommand{\El}{\mathsf{El}}
\newcommand{\Clock}{\mathsf{Clock}}

\newcommand{\Machine}[1]{\mathsf{Machine}\, {#1}}


% Predomains and EP pairs
\newcommand{\Nat}{\mathsf{Nat}}
\newcommand{\Dyn}{\mathsf{Dyn}}
\newcommand{\ty}[1]{\langle {#1} \rangle}
\newcommand{\li}{L_\mho}
\newcommand{\liclk}[1]{L_\mho [{#1}]}

\newcommand{\ext}[2]{\text{ext}\,{#1}\,{#2}}
\newcommand{\map}[2]{\text{map}\,{#1}\,{#2}}

\newcommand{\ltls}{\ltdyn}
\newcommand{\bisim}{\approx}
\newcommand{\semlt}{\le}
\newcommand{\semltbad}{\lesssim}

%\newcommand{\injarr}{\textsf{Inj}_\to}
%\newcommand{\injnat}{\textsf{Inj}_\mathbb{N}}

\newcommand{\id}{\mathsf{id}}
\newcommand{\ep}{\leadsto}

\newcommand{\emb}[2]{\mathsf{emb}_{#1}({#2})}
\newcommand{\proj}[2]{\mathsf{proj}_{#1}({#2})}

\newcommand{\monto}{\to_m}


% Notation for wait functions
\newcommand{\wre}{w_r^e}
\newcommand{\wle}{w_l^e}
\newcommand{\wrp}{w_r^p}
\newcommand{\wlp}{w_l^p}

\newcommand{\sem}[1]{\llbracket {#1} \rrbracket}
\newcommand{\semgl}[1]{{\sem{#1}}^\text{gl}}


% Denotational model
\newcommand{\errdom}{\textsf{ErrDom}}

\newcommand{\ptb}{\text{ptb}}
\newcommand{\push}{\text{push}}
\newcommand{\pull}{\text{pull}}

\newcommand{\pv}{P^{\mathcal{V}}}
\newcommand{\pe}{P^{\mathcal{E}}}
\newcommand{\ptbv}{\text{ptb}^\mathcal{V}}
\newcommand{\ptbe}{\text{ptb}^\mathcal{E}}

\newcommand{\Ppure}{P}
\newcommand{\Pk}{P^K}
\newcommand{\ptbk}{\text{ptb}^K}


\newcommand{\upf}{\text{up}}
\newcommand{\dnf}{\text{dn}
}
\newcommand{\arr}{\to}
\newcommand{\comp}{\bullet}

\newcommand{\ltsq}[2]{\mathrel{\ltdyn^{#1}_{#2}}}
\newcommand{\ltsqbisim}[2]{\mathrel{{\widetilde{\ltdyn}}^{#1}_{#2}}}
\newcommand{\ltbisim}{\mathrel{\widetilde{\ltdyn}}}
\newcommand{\vf}{\mathcal{V}_f}
\newcommand{\vsq}{\mathcal{V}_{sq}}
\newcommand{\ef}{\mathcal{E}_f}
\newcommand{\esq}{\mathcal{E}_{sq}}

\newcommand{\morbisimid}[1]{\text{Endo}_\bisim{#1}}


\newcommand{\sv}{s_{\mathcal{V}}}
\newcommand{\tv}{t_{\mathcal{V}}}
\newcommand{\rv}{r_{\mathcal{V}}}
\newcommand{\se}{s_{\mathcal{E}}}
\newcommand{\te}{t_{\mathcal{E}}}
\newcommand{\re}{r_{\mathcal{E}}}

\newcommand{\Ff}{F_f}
\newcommand{\Fsq}{F_{sq}}
\newcommand{\Uf}{U_f}
\newcommand{\Usq}{U_{sq}}
\newcommand{\arrf}{\arr_f}
\newcommand{\arrsq}{\arr_{sq}}

\newcommand{\vsim}{\mathcal{V}_{\bisim}}
\newcommand{\esim}{\mathcal{E}_{\bisim}}
\newcommand{\svsim}{s_{\mathcal{V}}^\bisim}
\newcommand{\tvsim}{t_{\mathcal{V}}^\bisim}
\newcommand{\rvsim}{r_{\mathcal{V}}^\bisim}
\newcommand{\sesim}{s_{\mathcal{E}}^\bisim}
\newcommand{\tesim}{t_{\mathcal{E}}^\bisim}
\newcommand{\resim}{r_{\mathcal{E}}^\bisim}



\newcommand{\ve}{\mathcal{V}_e}
\newcommand{\ee}{\mathcal{E}_e}

\newcommand{\vr}{\mathcal{V}_r}
\newcommand{\er}{\mathcal{E}_r}

\newcommand{\relatedin}[3]{{#1}\, {#3}\, {#2}}
\newcommand{\binrel}[1]{\mathbin{#1}}


\newcommand{\da}{\downarrow}

\newcommand{\ret}{\text{ret}}
\newcommand{\bind}[3]{{#1} <- {#2}\, ; \, {#3}}

\newcommand{\piv}{\Pi^\mathcal{V}}
\newcommand{\pie}{\Pi^\mathcal{E}}

\newcommand{\upl}{\textsc{UpL}}
\newcommand{\upr}{\textsc{UpR}}
\newcommand{\dnl}{\textsc{DnL}}
\newcommand{\dnr}{\textsc{DnR}}

\newcommand{\delre}{\delta^{r,e}}
\newcommand{\delle}{\delta^{l,e}}
\newcommand{\delrp}{\delta^{r,p}}
\newcommand{\dellp}{\delta^{l,p}}

\newcommand{\qordeq}{\bisim}

\newcommand{\inat}{\text{Inj}_\mathbb{N}}
\newcommand{\itimes}{\text{Inj}_\times}
\newcommand{\iarr}{\text{Inj}_\to}

\newcommand{\tnat}{\mathsf{nat}}
\newcommand{\ttimes}{\mathsf{times}}
\newcommand{\tfun}{\mathsf{fun}}

\newcommand{\cased}[7]{
    \begin{align*}
    \text{Case}_D ({#1}) &\text{ of } \{ \\
        &\alt \text{nat}({#2}) \to {#3}  \\
        &\alt \text{times}({#4}) \to {#5}  \\
        &\alt \text{fun}({#6}) \to {#7} \}
    \end{align*}
}

\newcommand{\inone}{\mathsf{in}_1}
\newcommand{\intwo}{\mathsf{in}_2}
\newcommand{\inthree}{\mathsf{in}_3}
\newcommand{\infour}{\mathsf{in}_4}
\newcommand{\infive}{\mathsf{in}_5}


\newcommand{\delay}{\text{Delay}}
\newcommand{\ledelay}{\le^{\text{Del}}}
\newcommand{\bisimdelay}{\bisim^{\text{Del}}}
\newcommand{\tnow}{\mathsf{now}}
\newcommand{\tlater}{\mathsf{later}}
\newcommand{\Da}{\Downarrow}




% Is there a notion of naturality for functors on double categories that
% involves horizontal morphisms?

% Come up with consistent notation for the double category of values of a double
% CBPV model, and likewise for computations




\begin{document}

\title{An Abstract, Guarded Categorical Model of Gradual Typing}
\author{Eric Giovannini}

\maketitle


The work of Giovannini et al. \cite{gtt-sgdt} constructs a denotational
semantics for a simple gradually-typed lambda calculus using guarded type
theory.
%
Many of the aspects of the model construction are in fact independent of the
details of the particular language being modeled, e.g., the specific types and
effects in the language.
%
In this work, we seek to abstract over these details and present an abstract
categorical framework for constructing models of gradual typing in guarded type
theory.

With this abstract framework in hand, we then define concrete semantics for
gradually-typed languages with simple effects.

\subsection{Why an Abstract Model?}

The model developed by Giovannini et al. suffers from the limitation that it is
a single model for a quite bare-bones gradually-typed language. Real-world
languages can include effects, polymorphism, dependent types, etc. The goals of
an abstract model are twofold:

\begin{enumerate}
  \item To promote a better understanding of the essence of any guarded
  relational semantics of gradual tying independent of the particular language
  features being modeled.
  \item To allow for developers of languages to establish graduality for their
  language by instantiating our abstract framework.
\end{enumerate}

\subsection{Challenges}

Prior work by New and Licata \cite{new-licata18} defines a categorical model for
gradual typing that serves as the basis for denotational semantics using
\emph{classical} domain theory. Their model is based on \emph{double categories}
\cite{shulman2007}.
%
Unfortunately, as Giovannini et al. point out, adapting this model to the
guarded setting poses significant challenges, and much of their work consists of
defining a version of the model that works in guarded type theory.

Inspired by the solution of Giovannini et al., we refine the notion of abstract
categorical model of gradual typing into two novel classes of model suitable for
the guarded setting. We call these \emph{extensional} and \emph{intensional}
models.

% TODO more explanation

\subsection{Outline}

This work is organized as follows.

\begin{itemize}
  \item We begin in Section \ref{sec:idealized-models} with an overview of the
  existing abstract double-categorical semantics for gradual typing by New and
  Licata.
  \item We review the difficulties in adapting the New-Licata denotational model
  to the guarded setting.
  \item We motivate the refinement of the notion of abstract model into
  extensional and intensional models and present their definitions in full.
  \item We discuss the construction of an extensional model of gradual typing
  from a suitable ``base'' double category of effectful functions and relations.
  \item With the constructions in hand, we then instantiate the abstract
  framework to define concrete models of gradual typing capable of modeling
  languages with simple effects.
\end{itemize}

\subsection{Ambient Guarded Metalanguage}

Because guarded type theory plays a central role in the construction of the
concrete model in prior work, it is natural to expect that our abstract
framework should to some degree incorporate aspects of guarded type theory.
%
To that end, we demand the following of the categories involved in the
description of our abstract model:

\begin{enumerate}
    \item There is an applicative functor $\later \colon \calC \to \calC$.
    \item There is a morphism $\nxt \colon X \to \later X$ for any object $X$.
    \item For any morphism $f$ of the form $\later X \to X$, there is a unique
    morphism $\fix f \colon 1 \to X$ such that $\fix f = f \circ \nxt \circ \fix
    f$.
\end{enumerate}


% \section{Extensional and Intensional Models}

\section{Double Categorical Semantics of Graduality}\label{sec:idealized-models}

We begin with a review of abstract categorical models for gradual typing
described in prior work.
%
New and Licata \cite{new-licata18} modeled the graduality and type casts for
call-by-name gradual typing using \emph{double categories}, which are defined to
be categories internal to the category of categories. That is, a double category
$\mathcal C$ consists of a category $\mathcal C_f$ of ``objects and function
morphisms'' and a category $\mathcal C_{sq}$ of ``relation morphisms and
squares'' with functors (reflexive relation) $r : C_f \to C_{sq}$ and (source
and target) $s,t : C_{sq} \to C_f$ satisfying $sr = tr = \id$ as well as a
composition operation $c : C_{sq} \times_{s,t} C_{sq} \to C_{sq}$ respecting
source and target. This models an abstract notion of functions and relations.
For notation, we write function morphisms as $f : A \to B$ and relation
morphisms as $c : A \rel B$ where $c \in C_{sq}$ and $s(c) = A$ and $t(c) = B$.
Finally, a morphism $\alpha$ from $c$ to $d$ with $s(\alpha) = f$ and $s(\beta)
= g$ is visualized as

% https://q.uiver.app/#q=WzAsNCxbMCwwLCJBIl0sWzEsMCwiQiJdLFswLDEsIkEnIl0sWzEsMSwiQiciXSxbMCwyLCJmIiwyXSxbMSwzLCJnIl0sWzAsMSwiYyIsMCx7InN0eWxlIjp7ImJvZHkiOnsibmFtZSI6ImJhcnJlZCJ9LCJoZWFkIjp7Im5hbWUiOiJub25lIn19fV0sWzIsMywiZCIsMix7InN0eWxlIjp7ImJvZHkiOnsibmFtZSI6ImJhcnJlZCJ9LCJoZWFkIjp7Im5hbWUiOiJub25lIn19fV1d
\[\begin{tikzcd}[ampersand replacement=\&]
	A \& B \\
	{A'} \& {B'}
	\arrow["f"', from=1-1, to=2-1]
	\arrow["g", from=1-2, to=2-2]
	\arrow["c", "\shortmid"{marking}, no head, from=1-1, to=1-2]
	\arrow["d"', "\shortmid"{marking}, no head, from=2-1, to=2-2]
\end{tikzcd}\] 

Such a morphism is thought of as an abstraction of the notion of relatedness of
functions: functions take related inputs to related outputs. The composition
operations and functoriality give us a notion of composition of relations as
well as functions and vertical and horizontal composition of squares. In this
work we will be chiefly interested in \emph{locally thin} double categories,
that is, double categories where there is at most one square for any $f,c,g,d$.
In this case we use the notation $f \ltsq{c}{d} g$ to mean that a square like the
above exists.

New, Licata and Ahmed \cite{new-licata-ahmed2019} extended the axiomatic syntax
to call-by-push-value but did not analyze the structure categorically. We fill
in this missing analysis now, extending the double categorical semantics given
in \cite{new-licata18} from cartesian closed categories to CBPV models. 

We begin by reviewing the definition of a CBPV model:

\begin{definition}
  A CBPV model $\calM$ consists of:
  \begin{itemize}
    \item A cartesian-closed category $\calV$ of values and pure morphisms.
    \item A category $\calE$ of computations and linear morphisms.
    \item Functors $F : \calV \to \calE$ and $U : \calE \to \calV$ with $F \dashv U$.
    \item An action of $\calV$ on $\calE$, i.e., a functor $\arr \colon
    \calV^{op} \times \calE \to \calE$ satisfying the usual coherence laws,
    and such that $U(A \arr B)$ is isomorphic to $A \Rightarrow UB$.
  \end{itemize}

  A \emph{strict homomorphism of CBPV models} from $(\mathcal V,\mathcal
  E,\ldots)$ to $(\mathcal V', \mathcal E',\ldots)$ consists of functors $G_v :
  \mathcal V \to \mathcal V'$ and $G_e : \mathcal E \to \mathcal E'$ that
  strictly preserve all CBPV constructions; see the appendix for a more detailed
  definition. We call this a strict morphism in contrast to a \emph{lax}
  morphism, which only preserves CBPV constructions up to transformation.
\end{definition}

A model of the congruence rules of the CBPV system can be given by a locally
thin ``double CBPV model'', which we define below:

\begin{definition}\label{def:double-cbpv-model}
  A \emph{double CBPV model} is a category internal to the category of CBPV
  models and \emph{strict} homomorphisms of CBPV models\footnote{it may be
  possible to also define this as a notion of CBPV model internal to some
  structured $2$-category of categories, but we are not aware of any such
  definition of an internal CBPV model}.
\end{definition}

Part of the data of a double CBPV model can be visualized as follows:

% https://q.uiver.app/#q=WzAsNCxbMCwwLCJcXHZzcSJdLFsyLDAsIlxcZXNxIl0sWzAsMiwiXFx2ZiJdLFsyLDIsIlxcZWYiXSxbMiwzLCJcXEZmIiwwLHsiY3VydmUiOi0yfV0sWzMsMiwiXFxVZiIsMCx7ImN1cnZlIjotMn1dLFswLDEsIlxcRnNxIiwwLHsiY3VydmUiOi0yfV0sWzEsMCwiXFxVc3EiLDAseyJjdXJ2ZSI6LTJ9XSxbMiwwLCJcXHJ2Il0sWzAsMiwiXFx0diIsMCx7Im9mZnNldCI6LTEsImN1cnZlIjotMn1dLFswLDIsIlxcc3YiLDIseyJvZmZzZXQiOjEsImN1cnZlIjoyfV0sWzEsMywiXFx0ZSIsMCx7Im9mZnNldCI6LTEsImN1cnZlIjotMn1dLFsxLDMsIlxcc2UiLDIseyJvZmZzZXQiOjEsImN1cnZlIjoyfV0sWzMsMSwiXFxyZSJdLFs0LDUsIlxcYm90IiwxLHsic2hvcnRlbiI6eyJzb3VyY2UiOjIwLCJ0YXJnZXQiOjIwfSwic3R5bGUiOnsiYm9keSI6eyJuYW1lIjoibm9uZSJ9LCJoZWFkIjp7Im5hbWUiOiJub25lIn19fV0sWzYsNywiXFxib3QiLDEseyJzaG9ydGVuIjp7InNvdXJjZSI6MjAsInRhcmdldCI6MjB9LCJzdHlsZSI6eyJib2R5Ijp7Im5hbWUiOiJub25lIn0sImhlYWQiOnsibmFtZSI6Im5vbmUifX19XV0=
\[\begin{tikzcd}[ampersand replacement=\&]
	\vsq \&\& \esq \\
	\\
	\vf \&\& \ef
	\arrow[""{name=0, anchor=center, inner sep=0}, "\Fsq", curve={height=-12pt}, from=1-1, to=1-3]
	\arrow["\tv", shift left, curve={height=-12pt}, from=1-1, to=3-1]
	\arrow["\sv"', shift right, curve={height=12pt}, from=1-1, to=3-1]
	\arrow[""{name=1, anchor=center, inner sep=0}, "\Usq", curve={height=-12pt}, from=1-3, to=1-1]
	\arrow["\te", shift left, curve={height=-12pt}, from=1-3, to=3-3]
	\arrow["\se"', shift right, curve={height=12pt}, from=1-3, to=3-3]
	\arrow["\rv", from=3-1, to=1-1]
	\arrow[""{name=2, anchor=center, inner sep=0}, "\Ff", curve={height=-12pt}, from=3-1, to=3-3]
	\arrow["\re", from=3-3, to=1-3]
	\arrow[""{name=3, anchor=center, inner sep=0}, "\Uf", curve={height=-12pt}, from=3-3, to=3-1]
	\arrow["\bot"{description}, draw=none, from=0, to=1]
	\arrow["\bot"{description}, draw=none, from=2, to=3]
\end{tikzcd}\]

The diagram above makes it clear that there are two useful ways of viewing a
double CBPV model $\calM$. First, we can view $\calM$ as pair of CBPV models,
one forming the objects and vertical morphisms (the bottom of the diagram) and
the other forming the horizontal morphisms and the squares (the top of the
diagram). Alternatively, we can view $\calM$ as a pair of double categories, one
for the values (the left of the diagram), and one for the computations (the
right of the diagram). For the latter viewpoint, we will often write $\calC$ to
stand for an unspecified such double category associated with a double CBPV
model.

Given a double CBPV model, we can interpret the syntax of gradual typing as
follows. Type precision $A \ltdyn A'$ is interpreted as a relation morphism $c_A
: A \rel A'$ in $\mathcal V_{sq}$, and term precision $\Gamma \ltdyn \Gamma'
\vdash M \ltdyn M' : A \ltdyn A'$ is interpreted as a square $M
\ltsq{c_\Gamma}{UF c_A} M'$.
%
% TODO review this
The fact that $t,r$ and the composition are all given by strict CBPV
homomorphisms says that all the type constructors lift to precision
(monotonicity of type constructors) as well as all term constructors
(congruence). Further, $r$ and composition being strict homomorphisms implies
that all type constructors strictly preserve the identity relation (identity
extension) and composition.

Next, to model type casts, the New-Licata model further requires that every
value relation $c : A \rel A'$ is \emph{left representable} by a function $u_c :
A \to A'$ and every computation relation $d : B \rel B'$ is \emph{right
representable} by a function $d_c : B' \to B$. In a locally thin double
category, left and right representability are defined as follows:

\begin{definition} % TODO these are the simpler versions of the rules
  A relation $R : A \rel B$ is \emph{left representable} by $f : A \to B$ if $f
  \ltsq{R}{r(B)} \id$ and $\id \ltsq{r(A)}{R} f $.

  Dually, $R : A \rel B$ is \emph{right representable} by $g : B \to A$ if $\id
  \ltsq{R}{r(A)} g$ and $g \ltsq{r(B)}{R} \id$.
\end{definition}

% These squares say that the relation $c$ is \emph{representable} by the upcast
% morphism $\upc{c}$, essentially that the relation is a kind of \emph{graph} of
% the function in that $x \binrel{\sem c} y$ if and only if $\sem{\upc c}(x)
% \leq_{A_2} y$ \cite{shulman2007}.

These squares are sufficient to model the UpL/UpR/DnL/DnR rules for casts.
% TODO this is possible because of horizontal composition of squares
Additionally, since representable morphisms compose, the compositionality of
casts comes for free. However, the retraction property must be added as an
additional axiom to the model. 
%
To model the error being the least element, the New-Licata model adds the
requirement that $\mho \circ {!} \ltsq{r(A)}{r(UB)} f$ holds for all $f :
\mathcal V(A,B)$. Finally, the dynamic type can be modeled as an arbitrary value
type $D$ with arbitrary relations $\nat \rel D$ and $D \ra D \rel D$ and $D
\times D \rel D$ (or whatever basic type cases are required).


\section{Adapting the Notion of Model to Guarded Type Theory}


% we have seen in the previous section that there are obstructions to developing a
% concrete model in \emph{guarded} type theory. While guarded type theory makes
% the construction of arbitrary guarded recursive definitions possible, it comes
% at a significant cost to reasoning: unfolding of recursive definitions is
% explicit. In a denotational semantics, this means that non-well-founded
% recursion must perform a kind of \emph{observable step}. This presents a
% difficulty in proving graduality, which is a property that is oblivious to the
% number of steps that a program takes.

% Therefore, to prove graduality, we must be able to relate terms that take
% different numbers of steps. In the previous section we have examined two ways of
% doing this: in the first approach by defining an ordering $\semltbad$ that is
% closed under delays, and in the second by decomposing this ordering into a
% lock-step error ordering and a weak bisimilarity relation. But either way, we
% cannot escape the confines of Theorem \ref{thm:no-go}, which tells us that
% neither of these fomulations is transitive.

\subsection{Challenges in Adapting the New-Licata Model}

While the original double categorical semantics introduced in Section
\ref{sec:idealized-models} can be satisfied with classical domain theoretic
models as demonstrated by New and Licata \cite{new-licata18}, Giovannini et al.
\cite{gtt-sgdt} show that the New-Licata model cannot be directly adapted to
work in the guarded setting.

Key to understanding the issue is the fact that the equation for the dynamic
type involves a later, which means that when inspecting an element of the
dynamic type which is a function, an \emph{observable step} is required in the
semantics. This means that downcasts from the dynamic type may take an
observable step, and thus \emph{any} cast may take steps as it may involve a
downcast from the dynamic type functorially.

The issue in adapting the New-Licata model to the guarded setting is related to
the representability condition on the relations. At a high level, the problem
arises as a result of a tension between \emph{transitive reasoning} and
\emph{step-insensitivity}.
%
Transitive reasoning is desirable because it allows us to derive the cast rules
needed for graduality from the simpler, compositional versions of the rules.
This is crucial because the original rules quantify over an additional,
arbitrary type precision derivation and are therefore not syntax-independent.
%
On the other hand, the graduality property, as axiomatized by the
representability of type precision relations, is fundamentally a
step-insensitive property. The casts, which may involve observable steps, must
be related to the identity function, which does not step. The relation
modeling term precision must be able to account for this mismatch.

However, Giovannini et al. show that transitivity and unrestricted
step-insensitivity (plus a third minor condition) are incompatible as properties
of a relation in guarded type theory (see their Theorem 4.1 for a precise
statement).
%
The solution adopted by Giovannini et al. is to work with a \emph{lock-step
error ordering} that is transitive, and then restore a small amount of
step-insensitivity where needed through explicit synchronizing terms they call
\emph{perturbations}. The perturbations mimic the stepping behavior of casts,
but are otherwise equivalent to the identity, a property they formalize via a
notion of \emph{weak bisimilarity}. The authors relax the definition of
representability to allow for casts to be lock-step related to perturbations
rather than only to the identity function.
%
Finally, to define the denotation of term precision, the authors take the
closure of the lock-step ordering under weak bisimilarity on both sides:
%
\[ M \ltbisim N := \exists f, g.\, M \bisim f \ltdyn g \bisim N, \] 
%
where $\bisim$ is weak bisimilarity and $\ltdyn$ is the lock-step ordering.
%
The idea is that two terms are related if they can be synchronized in some
number of steps to two terms that are \emph{lock-step} related.
% This definition allows them to relate terms that differ up to a perturbation. 
% 
By separating term precision into a combination of weak bisimilarity and
lock-step reasoning, and using syntactic perturbations as a means to explicitly
synchronize terms as necessary, the authors achieve compositional reasoning
while still accommodating the steps taken by casts.




\subsection{Intensional Models of Gradual Typing}

An intensional model of gradual typing is an abstraction of working in a setting
where the ordering on terms is transitive but requires terms to be in lock-step.
To aid in understanding the definition of an intensional model, we break the
definition into five steps, with each one building on the previous. The last
step will be the definition we have in mind when we refer to an intensional
model.

% The steps in the definition of intensional model are as follows:
% \begin{enumerate}
%   \item A double CBPV model
% \end{enumerate}

The starting point for an intensional model is the same as that of an idealized
model of gradual typing: namely, a double CBPV model (see Definition
\ref{def:double-cbpv-model}). In particular, unlike the extensional model where
we could only compose \emph{relations} horizontally, here we also have a
horizontal composition operation on squares.


As in the definition of extensional model, we also require the existence of a
natural transformation $\mho : 1 \Rightarrow U$ such that $\mho\, \circ\,{!}
\ltsq{r(A)}{r(UB)} f$ for any $f : A \to UB$, and a distinguished value type
$\nat$ with morphisms $z : \mathcal V(1,\nat)$ and $s : \mathcal V(\nat,\nat)$.

We next discuss the aspects that allow us to model intensional behavior. To
begin, we need guarded primitives.

\begin{enumerate}
    \item There is an applicative endofunctor $\later \colon \calC \to \calC$.
    \item There is a natural transformation $\nxt \colon \id \to \later$.
    \item For any morphism $f$ of the form $\later X \to X$, there is a unique
    morphism $\fix f \colon 1 \to X$ such that $\fix f = f \circ \nxt \circ \fix
    f$.
\end{enumerate}

% TODO verify
Notice that by naturality, we have that $(\later g) \circ \nxt = \nxt \circ g$
for any $g \colon X \to Y$.

We further require there to be a distinguished ``think'' map $\theta_{B} \colon
\later(UB) \to UB$ for all computation types $B$. From this, we can define the
``delay'' map $\delta_B := \theta_B \circ \nxt : UB \to UB$.
%
Moreover, we require that the think maps are related in that for any $d : B \rel
B'$, we have a square $\theta_{B} \ltsq{\later (Ud)}{Ud} \theta_{B'}$. We also
require that these maps commute with computation morphisms, in the sense that
for any $\phi \in \ef(B_1, B_2)$ we have $U\phi \circ \theta_{B_1} = \theta_{B_2}
\circ \later(U\phi)$.

We call a model satisfying these requirements an \emph{intensional pre-model}.
For the sake of ease of reference, we recap the definition below:
%
\begin{definition}
  An \textbf{intensional pre-model} consists of:
  \begin{itemize}
    \item A double CBPV model, i.e., a category internal to the category of CBPV
    models and lax morphisms, where we require that the morphisms $r$, $s$, and
    $t$ are strict.
    \item A natural transformation $\mho : 1 \Rightarrow U$ such that $\mho
    \circ {!} \ltsq{r(A)}{r(UB)} f$ for any $f : A \to UB$. 
    \item A value type $\nat$ with morphisms $z : \mathcal V(1,\nat)$ and $s :
    \mathcal V(\nat,\nat)$.
    \item A distinguished ``think'' map $\theta_{B} \in \vf(\later(UB), UB)$
    for all computation types $B$. We define $\delta_B := \theta_B \circ \nxt
    \colon \vf(UB, UB)$.
    \item For any $d : B \rel B'$, there is a square $\theta_{B}
    \ltsq{\later(Ud)}{Ud} \theta_{B'}$. 
    \item For any $\phi \in \ef(B, B')$ we have $U\phi \circ \theta_{B_1} =
    \theta_{B_2} \circ \later(U\phi)$.
  \end{itemize}
\end{definition}


\begin{lemma}
  We have the following properties involving the delay maps in any intensional
  pre-model:
  \begin{enumerate}
    \item For any $d : B \rel B'$, we have a square $\delta_{B} \ltdyn_{Ud}^{Ud} \delta_{B'}$
    \item We have $U\phi \circ \delta_{B_1} = \delta_{B_2} \circ U\phi$ for
    any $\phi \in \ef(B, B')$.
    \item We can define a computation morphism $\delta_{FA}^* \in \ef(FA, FA)$ for
    all $A$. To do so, first compose the unit $\eta_A \in \vf(A, UFA)$ with the
    morphism $\delta_{FA} \in \vf(UFA, UFA)$, and then by the adjunction we get a
    computation morphism $\delta_{FA}^* \in \ef(FA, FA)$.
    \item There is a square $\delta_{FA}^* \ltdyn_{Fc}^{Fc} \delta_{FA'}^*$ for all $c : A \rel A$.
  \end{enumerate}
\end{lemma}
\begin{proof}
  See appendix.
\end{proof}

%   U\phi \circ \delta_{B_1}
% = U\phi \circ \theta_{B_1} \circ \nxt 
% = \theta_{B_2} \circ \later(U\phi) \circ \nxt
% = \theta_{B_2} \circ \nxt \circ U\phi
% = \delta_{B_2} \circ U\phi

\subsubsection{Weak Bisimilarity}\label{sec:abstract-model-bisimilarity}

The next ingredient in the definition of intensional model is \emph{weak
bisimilarity}, which formalizes the notion that two morphisms are the same ``up
to steps'', i.e., they differ only in that one may wait more than the other.
%
In particular, for any pair of objects $A$ and $A'$, in $\vf$, we require that
there is a reflexive, symmetric relation $\bisim_{A,A'}$ on the hom-set $\vf(A,
A')$, called the \emph{weak bisimilarity} relation. Similarly for the
computation category: there is a reflexive, symmetric relation $\bisim_{B,B'}$
defined on each hom-set $\ef(B, B')$.
%
Additionally, the weak bisimilarity relation should be a congruence with respect
to composition: if $f \bisim_{A,A'} f'$ and $g \bisim_{A',A''} g'$, then $g
\circ f \bisim_{A,A''} g' \circ f'$, and likewise for computations.

The bisimilarity relation should also be preserved by the functors $U, F, \arr$,
and $\times$ e.g., if $f \bisim_{A',A} g$ and $\phi \bisim_{B,B'} \phi'$, then
$f \arr \phi \bisim_{A \arr B, A' \arr B'} g \arr \phi'$. Similarly, the two
directions of the isomorphism of hom-sets $\vf(A, UB) \cong \ef(FA, B)$ should
preserve bisimilarity.

We can state all of this compactly by requiring that the categories $\vf$ and
$\ef$ be \emph{enriched} over the symmetric monoidal category $\setapprox$,
whose objects are sets equipped with a reflexive, symmetric relation and whose
arrows are functions between the sets that preserve the relation. Then we
require that the functors $U, F, \arr$, and $\times$ are
$\setapprox$-enriched, and likewise $F \dashv U$ as enriched functors.
%
Note that $\vsq$ and $\esq$ are not $\setapprox$-enriched, so in the functors
$r, s, t$ between $\vf$ and $\vsq$ and between $\ef$ and $\esq$, we view $\vf$
and $\ef$ as ordinary categories, where we simply forget the bisimilarity
relation on hom-sets.

% TODO do we need to require this?
% It seems this is needed when defining the actions of the functors on semantic perturbations.
%
% The weak bisimilarity relation should furthermore be preserved by the functors
% $U, F, \arr,$ and $\times$, e.g., if $f \bisim_{A',A} g$ and $\phi \bisim_{B,B'} \phi'$, then
% $f \arr \phi \bisim_{A \arr B, A' \arr B'} g \arr \phi'$.

Lastly, we require that for any computation object $B$, the delay map $\delta_B
\colon UB \to UB$ (which is defined in terms of $\theta_B$) is bisimilar to the
identity: $\delta_B \bisim_{UB, UB} \id_{UB}$.


\begin{definition}\label{def:step-1-model}   
An \emph{intensional pre-model with bisimilarity} is an intensional pre-model
$\calM$ such that:
\begin{itemize}
  % \item There is a reflexive, symmetric relation $\bisim$ on the hom-sets
  % $\vf(A,A')$ and $\ef(B,B')$ such that $\bisim$ preserves composition.
  \item The categories $\vf$ and $\ef$ are enriched over $\setapprox$.
  \item The functors $U, F, \arr$, and $\times$ are $\setapprox$-enriched, and
  the adjunction $F \dashv U$ is an enriched adjunction.
  \item For each $B$, the delay map $\delta_{B}$ is bisimilar to the identity
  $\delta_B \bisim_{UB, UB} \id_{UB}$.
\end{itemize}
\end{definition}

If $\calM$ is an intensional pre-model with bisimilarity, it will sometimes be
useful to have a concise definition describing the value objects, morphisms,
relations, and squares as a coherent unit, and likewise for the computations. To
do this, we introduce some terminology:

\begin{definition}
  A \emph{bisimilarity reflexive-graph} $\calC$ consists of:
  \begin{itemize}
    \item A category $\cf$ of objects and vertical morphisms, enriched in
    $\setapprox$.
    \item A thin category $\csq$ of horizontal morphisms and squares.
    \item A functor $r: |\cf| \to \csq$ where $|\cf|$ is the underlying ordinary
    category of $\cf$.
    \item Functors $s, t: \csq \to |\cf|$ such that $sr = tr = \id$.
  \end{itemize}

  A \emph{functor} $G : \calC \to \calC'$ of bisimilarity reflexive-graphs
  consists of:
  \begin{itemize}
    \item A $\setapprox$-enriched functor $G_f: \cf \to \cf'$.
    \item A functor $G_{sq} : \csq \to \csq'$.
    \item Source and target are preserved: $s' \circ G_{sq} = G_f \circ s$, and $t' \circ G_{sq} = G_f \circ t$
    \item Reflexivity is preserved: $r' \circ G_f = G_{sq} \circ r$.
  \end{itemize}

  Let $G, H: \calC \to \calC'$ be functors of bisimilarity reflexive-graphs. A
  \emph{natural transformation} $\alpha: G \to H$ consists of:
  \begin{itemize}
    \item A $\setapprox$-enriched natural transformation $\alpha^f: G_f \to H_f$.
    \item A natural transformation $\alpha^{sq}: G_{sq} \to H_{sq}$.
  \end{itemize}

  A \emph{bisimilarity double-category} is a bisimilarity reflexive graph that
  additionally has a functor $\circ \colon \csq \times_{s = t} \csq \to \csq$
  where $\csq \times_{s = t} \csq$ is the following pullback:

    % https://q.uiver.app/#q=WzAsNCxbMCwwLCJcXGNzcSBcXHRpbWVzX3tzID0gdH0gXFxjc3EiXSxbMCwxLCJcXGNzcSJdLFsxLDAsIlxcY3NxIl0sWzEsMSwiXFxjZiJdLFswLDFdLFsxLDNdLFsyLDNdLFswLDJdLFswLDMsIiIsMSx7Im9mZnNldCI6Miwic3R5bGUiOnsibmFtZSI6ImNvcm5lciJ9fV1d
  \[\begin{tikzcd}[ampersand replacement=\&]
    {\csq \times_{s = t} \csq} \& \csq \\
    \csq \& \cf
    \arrow[from=1-1, to=1-2]
    \arrow[from=1-1, to=2-1]
    \arrow["\lrcorner"{anchor=center, pos=0.125}, shift right=2, draw=none, from=1-1, to=2-2]
    \arrow[from=1-2, to=2-2]
    \arrow[from=2-1, to=2-2]
  \end{tikzcd}\]

  We additionally require that composition behaves with respect to the source and
  target functors, and satisfies associativity and left and right unit laws, as in
  an ordinary double category.

% % https://q.uiver.app/#q=WzAsNSxbMSwwLCJcXGNzcSBcXHRpbWVzX3tzID0gdH0gXFxjc3EiXSxbMSwyLCJcXGNmIl0sWzEsMSwiXFxjc3EiXSxbMCwxLCJcXGNzcSJdLFsyLDEsIlxcY3NxIl0sWzAsMiwiXFxjaXJjIl0sWzIsMSwicyIsMix7Im9mZnNldCI6MX1dLFszLDEsInMiLDJdLFswLDMsIlxccGlfMSIsMl0sWzAsNCwiXFxwaV8yIl0sWzQsMSwidCJdLFsyLDEsInQiLDAseyJvZmZzZXQiOi0xfV1d
% \[\begin{tikzcd}[ampersand replacement=\&]
% 	\& {\csq \times_{s = t} \csq} \\
% 	\csq \& \csq \& \csq \\
% 	\& \cf
% 	\arrow["{\pi_1}"', from=1-2, to=2-1]
% 	\arrow["\circ", from=1-2, to=2-2]
% 	\arrow["{\pi_2}", from=1-2, to=2-3]
% 	\arrow["s"', from=2-1, to=3-2]
% 	\arrow["s"', shift right, from=2-2, to=3-2]
% 	\arrow["t", shift left, from=2-2, to=3-2]
% 	\arrow["t", from=2-3, to=3-2]
% \end{tikzcd}\]


\end{definition}

With these definitions in hand, we can give an alternative definition of an
intensional pre-model with bisimilarity:

\begin{definition}

\end{definition}


\subsubsection{Perturbations}\label{sec:abstract-model-perturbations}

We next focus on the specification of perturbations in a pre-model with
bisimilarity. To begin, let $\calC$ be a bisimilarity double-category.
%
% Let $\calM$ be a pre-model with bisimilarity, and let $\calC$ stand
% for the double category of values of $\calM$, or the double category of
% computations of $\calM$.
%
For an object $X$ of $\calC$, a \emph{semantic perturbation} of $X$ is an
endomorphism $f \in \cf(X, X)$ such that $f \bisim_{X,X} \id_X$. The collection
of semantic perturbations on $X$ forms a monoid under composition, with the
identity element being the identity $\id_X$. We write $\semptb{X}$ to refer to
this monoid.

For every object $X \in \cf$, we require that there is a commutative
monoid\footnote{The model of Giovannini et al. used (not necessarily
commutative) monoids, whereas we demand that the monoids be commutative. The
reason for this will become clear in the next section.} $M_X$ of \emph{syntactic
perturbations} of $X$ and a monoid homomorphism $i_X : M_X \to \semptb{X}$ that
interprets these as semantic perturbations.

The prototypical semantic perturbation for a computation object $B$ is the delay
map $\delta_B \colon UB \to UB$. We require that for any computation object $B$,
the monoid $M_{UB}$ contains an element $m$ such that $i_{UB}(m) = \delta_B$.

We require that the functors $\times, \timesk$, $\arr, \tok$, $U, \Uk$, $F, \Fk$
all preserve perturbations. As in the prior work of Giovannini et al., we also
require that the perturbations interact nicely with the relation morphisms.
Namely, if $R : X \rel X'$ is a relation morphism, we require there to be two
homomorphisms of monoids $\push_R : M_X \to M_{X'}$ and $\pull_R : M_{X'} \to
M_X$. The authors call this the \emph{push-pull property}.

If $\calM$ is a pre-model with bisimilarity that additionally satisfies the above
properties, we say that $M$ is a \emph{pre-model with perturbations}. To
summarize:
%
% TODO
\begin{definition}\label{def:step-2-model} 
  A \emph{pre-model with perturbations} $\calM$ is a pre-model with
  bisimilarity, such that the following additional requirements are met:
  \begin{enumerate}
    \item For each object (value or computation) $X$, there is a commutative
    monoid $M_X$ and a homomorphism of monoids $i_X : M_X \to \semptb{X}$.
    \item For all $B$, there is an element $m$ such that $i_{UB}(m) = \delta_B$.
    \item The functors $\times, \timesk$, $\arr, \tok$, $U, \Uk$, $F, \Fk$ preserve perturbations.
    \item The push-pull property holds for all $c : A \rel A'$ and all $d : B \rel B'$.
  \end{enumerate}
\end{definition}


\subsection{Representable Relations}\label{sec:abstract-model-representable}

Now we turn to the notion of representability of relations. Here we make a
departure from prior work: rather than working with quasi-representability, we
instead work with \emph{actual} representability in a certain double category of
\emph{well-behaved morphisms}. The idea is as follows.
%
Recall that in the model of Giovannini et al., the vertical morphisms had no
interaction with the monoid of perturbations; using their terminology, a
morphism of value objects (predomains equipped with syntactic perturbations) was
simply a morphism of the underlying predomains.
%
We observe, however, that many morphisms actually \emph{do} interact nicely with
perturbations. In particular, the cast morphisms enjoy this property. (The
precise sense in which we mean ``interact nicely'' will be explained below.)
This observation leads to the following idea. If we define a notion of square
where the sides build in a perturbation, we can look at representable relations
in this double category rather than defining a separate, ad-hoc notion of
quasi-representable relations. The interaction between morphisms and
perturbations turns out to be precisely what we need for these squares to
compose, i.e., for these squares to form a double category.

More concretely, given a pre-model with perturbations $\calM$, we will define a
double category $\calV^*$ whose vertical morphisms consist of value morphisms
that behave well with respect to casts, and whose squares are squares in which
both sides are composed with perturbations. Likewise, we can define analogously
a double category $\calE^*$.
% 
In particular, let $\calC$ stand for the double category of values of $\calM$,
or the double category of computations of $\calM$. We first define what we mean
for a vertical morphism to have an action on perturbations:

\begin{definition}\label{def:commutes-with-perturbations}
  Let $f \in \calC(X_i, X_o)$. We say that $f$ \emph{commutes with
  perturbations} if:
  \begin{itemize}
    \item There are monoid homomorphisms 
      $\pushmor{f} : M_{X_i} \to M_{X_o}$ and 
      $\pullmor{f} : M_{X_o} \to M_{X_i}$.
    \item For all $m_i \in M_{X_i}$ and $m_o \in M_{X_o}$ we have 
      \[ (i(m_o) \circ f) \ltsq{r(X_i)}{r(X_o)} (f \circ i(\pullmor{f}(m_o))) \]
      and
      \[ (i(\pushmor{f}(m_i)) \circ f) \ltsq{r(X_i)}{r(X_o)} (f \circ i(m_i)). \]
  \end{itemize}
\end{definition}

\begin{remark}
  Not all vertical morphisms in $\calM$ commute with perturbations.
  % TODO finish
\end{remark}

\begin{remark}
  The reason for ordering rather than equality is that semantic perturbations
  for computation types are computation morphisms, and hence need to commute
  with the algebraic structure. Therefore, performing the perturbation after
  performing a computation morphism may not result in any delay being inserted,
  while performing the perturbation first will always insert a delay.
  %
  At first glance this would seem to be problematic, since one side has a delay
  where the other does not. However, the key idea is that the vertical morphisms
  are meant to be casts, which have no effects other than to delay or to error.
  In either case, we will be able to show that the side that performs the
  perturbation after the morphism is \emph{below} the side that performs the
  perturbation first.
\end{remark}

A square in this double category will be defined as a square in the usual sense,
but where we are allowed to synchronize the vertical morphisms on either side
with a perturbation. More concretely:

\begin{definition}\label{def:perturbed-square} Let $f : X_i \to X_o$ and $g :
X_i' \to X_o'$ be perturbation-preserving morphisms and let $R_i : X_i \rel
X_i'$ and $R_o : X_o \rel X_o'$ be horizontal morphisms.
%
We say that there is a \textbf{perturbed square} between $f$ and $g$, written $f
\ltsqdelay{R_i}{R_o} g$, if there \emph{merely exist} perturbations $m_f, m_f',
m_g, m_g'$ such that 
$i(m_f') \circ f \circ i(m_f) \ltsq{R_i}{R_o} i(m_g') \circ g \circ i(m_g)$.
%

\begin{center}
  \begin{tabular}{ m{7em}  m{1cm} m{7em} } 
    % https://q.uiver.app/#q=WzAsNCxbMCwwLCJYX2kiXSxbMCwxLCJYX28iXSxbMSwwLCJYX2knIl0sWzEsMSwiWF9vJyJdLFswLDEsImYiLDJdLFsyLDMsImciXSxbMCwyLCJSX2kiLDAseyJzdHlsZSI6eyJib2R5Ijp7Im5hbWUiOiJiYXJyZWQifSwiaGVhZCI6eyJuYW1lIjoibm9uZSJ9fX1dLFsxLDMsIlJfbyIsMix7InN0eWxlIjp7ImJvZHkiOnsibmFtZSI6ImJhcnJlZCJ9LCJoZWFkIjp7Im5hbWUiOiJub25lIn19fV1d
    \[\begin{tikzcd}[ampersand replacement=\&]
      {X_i} \& {X_i'} \\
      {X_o} \& {X_o'}
      \arrow["{R_i}", "\shortmid"{marking}, no head, from=1-1, to=1-2]
      \arrow["f"', from=1-1, to=2-1]
      \arrow["g", from=1-2, to=2-2]
      \arrow["{R_o}"', "\shortmid"{marking}, no head, from=2-1, to=2-2]
    \end{tikzcd}\] &
    %
    $\leadsto$ &
    %
    % https://q.uiver.app/#q=WzAsOCxbMCwwLCJYX2kiXSxbMCwxLCJYX2kiXSxbMCwyLCJYX28iXSxbMSwwLCJYX2knIl0sWzAsMywiWF9vIl0sWzEsMSwiWF9pJyJdLFsxLDIsIlhfbyciXSxbMSwzLCJYX28nIl0sWzAsMSwiaShtX2YpIiwyXSxbMSwyLCJmIiwyXSxbMiw0LCJpKG1fZicpIiwyXSxbMyw1LCJpKG1fZykiXSxbNSw2LCJnIl0sWzYsNywiaShtX2cnKSJdLFswLDMsIlJfaSIsMCx7InN0eWxlIjp7ImJvZHkiOnsibmFtZSI6ImJhcnJlZCJ9LCJoZWFkIjp7Im5hbWUiOiJub25lIn19fV0sWzQsNywiUl9vIiwyLHsic3R5bGUiOnsiYm9keSI6eyJuYW1lIjoiYmFycmVkIn0sImhlYWQiOnsibmFtZSI6Im5vbmUifX19XV0=
    \[\begin{tikzcd}[ampersand replacement=\&]
      {X_i} \& {X_i'} \\
      {X_i} \& {X_i'} \\
      {X_o} \& {X_o'} \\
      {X_o} \& {X_o'}
      \arrow["{R_i}", "\shortmid"{marking}, no head, from=1-1, to=1-2]
      \arrow["{i(m_f)}"', from=1-1, to=2-1]
      \arrow["{i(m_g)}", from=1-2, to=2-2]
      \arrow["f"', from=2-1, to=3-1]
      \arrow["g", from=2-2, to=3-2]
      \arrow["{i(m_f')}"', from=3-1, to=4-1]
      \arrow["{i(m_g')}", from=3-2, to=4-2]
      \arrow["{R_o}"', "\shortmid"{marking}, no head, from=4-1, to=4-2]
    \end{tikzcd}\] \\ 
  \end{tabular}
\end{center}

% We call such a square a \emph{perturbed square} to distinguish it from an ordinary square.
\end{definition}

Perturbed squares can be composed vertically and horizontally, as shown in
the following two lemmas: % See the Appendix for proofs.

\begin{lemma}
  Suppose $f_1 \ltsqdelay{R_1}{R_2} g_1$ and $f_2 \ltsqdelay{R_2}{R_3} g_2$.
  Then $(f_2 \circ f_1) \ltsqdelay{R_1}{R_3} (g_2 \circ g_1)$.
\end{lemma}

\begin{lemma}
  If $f \ltsqdelay{R_i}{R_o} g$ and $g \ltsqdelay{R_i'}{R_o'} h$, then
  $f \ltsqdelay{R_i \odot R_i'}{R_o \odot R_o'} h$.
\end{lemma}

Vertical composition uses the fact that the vertical morphisms preserve
perturbations, while the horizontal composition uses the fact that the monoids
are commutative, as well as the push-pull property of the relations. For
details, see the proofs in the appendix.

With these definitions, we define a locally-thin double category $\calC^*$ of
morphisms commuting with perturbations:
\begin{definition}\label{def:category-of-ptb-commuting-morphisms}
  Let $\calM$ be an intensional model with perturbations, and let $\calC$ be a
  bisimilarity double-category. We define a new double category $\calC^*$ as
  follows:
  \begin{itemize}
    \item The objects of $\calC^*$ are the same as those of $\calC$.
    \item The vertical morphisms of $\calC^*$ are those morphisms of $\calC$ that
    commute with perturbations.
    \item The horizontal morphisms of $\calC^*$ are those of $\calC$.
    \item The squares of $\calC^*$ are perturbed squares.
  \end{itemize}
\end{definition}
%
Note that this is locally-thin because the perturbations in the definition of
perturbed square are existentially quantified.
%
We can now define a notion of an intensional model where the horizontal
morphisms are representable:
\begin{definition}\label{def:step-3-model}
  Let $\calM$ be an intensional pre-model with perturbations. We say that
  $\calM$ is a \emph{pre-model with representability} if:
  \begin{itemize}
    \item Every value horizontal morphism $c : A \rel A'$ is left-representable
    in the double category $\calV^*$. That is, there is a
    perturbation-preserving morphism $\upc{c} \in \calV^*(A, A')$ that
    left-represents $c$.
    \item Every computation horizontal morphism $d : B \rel B'$ is
    right-representable in $\calE^*$. That is, there is a
    perturbation-preserving morphism $\dnc{d} \in \calE^*(B', B)$ that
    right-represents $d$.
  \end{itemize}
\end{definition}

\subsection{The Dynamic Type}\label{sec:abstract-model-dyn}

We now discuss how to specify the denotation of the dynamic type in an
intensional model. It is important to note that it makes sense to axiomatize the
dynamic type already in a base intensional pre-mdoel. That is, we do not require
the notions of weak bisimilarity, perturbations, or representability of
horizontal morphisms to define the denotation of the dynamic type.
% However, for our purposes, we will assume we are given an intensional
% pre-model with representability.

% Given such a model $\calM$, to interpret the dynamic type we require that
% there is...

Let $\calM$ be an intensional pre-model, a pre-model with bisimilarity, a
pre-model with perturbations, or a pre-model with representability. To interpret
the dynamic type, we require that there is a distinguished value object $D \in
\calV$, such that for every value object $A$, there is a distinguished
horizontal morphism $\inj{A} \colon A \rel D$.

Now we can state the definition of an intensional model:
\begin{definition}
  An \emph{intensional model of gradual typing} is an intensional pre-model with
  representability that furthermore has an interpretation of the dynamic type.
\end{definition}

\section{Constructing a Guarded Model of Gradual Typing}

In this section, we discuss how to \emph{construct} an abstract model of gradual
typing via a step-by-step process. Recall that the ultimate goal is to build an
extensional model of gradual typing (Definition \ref{def:extensional-model}), as
such a model contains the structure necessary to interpret the syntax and
precision rules of a gradual cast calculus.

Although the definition of an extensional model involves a significant amount of
data, it is not necessary to create an instance of an extensional model from
scratch. Much of the difficult work can be carried out abstractly. In
particular, it is possible to construct an extensional model from an intensional
model using the abstract analogue of the bisimilarity-closure construction in
the work of Giovannini et al. Moreover, much of the construction of an
intensional model can also be carried out abstractly.

\subsection{Starting Point: Intensional Pre-Model with Bisimilarity}

We take as a starting point an intensional pre-model with bisimilarity
(Definition \ref{def:step-1-model}).

\subsection{Adding Perturbations}

Suppose we are given an intensional pre-model with bisimilarity $\calM$. Our
goal is to construct a model $\calM'$ with perturbations, as in Definition
\ref{def:step-2-model}. Furthermore, the vertical morphisms in $\calM'$ will
commute with perturbations, as defined in Definition
\ref{def:commutes-with-perturbations}.

The summary of the construction is as follows: a value object in $\calM'$
consists of a value object $A$ of $\calM$ paired with a commutative monoid
$\Ptb_A$ of syntactic perturbations, and a homomorphism of monoids $\ptb_A :
\Ptb_A \to \semptb{A}$ interpreting the perturbations as semantic perturbations.
For reasons that will become clear when we define the functor $F$, we also add a
second commutative monoid of \emph{impure} or \emph{Kleisli} perturbations
$\Ptbk_A$ along with a monoid homomorphism $\ptbk_A : \Ptbk_A \to \semptb{FA}$.
%
Dually, a computation object of $\calM'$ consists of a computation object $B$ of
$\calM$ paired with commutative monoids $\Ptb_B$ and $\Ptbk_B$, along with
monoid homomorphisms $\ptb_B : \Ptb_B \to \semptb{B}$ and $\ptbk_B : \Ptbk_B \to
\semptb{UB}$.

A value morphism $f \in \vf'(A_i \to A_o)$ in $\calM'$ is defined to be a value
morphism in $\calM$ such that $f$ commutes with the perturbations in
$\Ptb_{A_i}$ and $\Ptb_{A_o}$, and $Ff$ commutes with the perturbations in
$\Ptbk_{A_i}$ and $\Ptbk_{A_o}$.
%
Dually, a computation morphism $\phi \in \ef'(B_i, B_o)$ is defined to be a
computation morphism in $\calM$ that commutes with the perturbations in
$\Ptb_{B_i}$ and $\Ptb_{B_o}$, and such that $U\phi$ commutes with the
perturbations in $\Ptbk_{B_i}$ and $\Ptbk_{B_o}$.

A value relation $c: A \rel A'$ in $\calM'$ is defined to be a value relation in
$\calM$ that satisfies the push-pull property with respect to the monoids
$\Ptb_A$ and $\Ptb_{A'}$, and such that $Fc$ satisfies push-pull with respect to
$\Ptbk_A$ and $\Ptbk_{A'}$.
%
Dually, a computation relation $d: B \rel B'$ in $\calM'$ is defined to be a
computation relation in $\calM$ that satisfies the push-pull property with
respect to $\Ptb_B$ and $\Ptb_{B'}$, and such that $Ud$ satisfies push-pull with
respect to $\Ptbk_B$ and $\Ptbk_{B'}$.

Value and computation squares in $\calM'$ are defined to be the same as those in
$\calM$.

Next, recall from Section \ref{TODO} that we can construct a bisimilarity
double-category whose vertical morphisms commute with perturbations and whose
squares involve pre- and post-composing each vertical morphism with a
perturbation.

\subsection{Adding Representability}

\subsection{The Dynamic Type}



\subsection{Extracting an Extensional Model: The Bisimilarity-Closure Construction}

Lastly, we can obtain an extensional model from an intensional model using the
abstract analogue of the bisimilarity-closure definition in the work of
Giovannini et al. In particular, let $\calM$ be an intensional model. We define
an extensional model $\calM'$ as follows. The objects and vertical/horizontal
morphisms are the same as those in $\calM$; it is only the definition of squares
that differs.

Specifically, let $f : X_i \to X_o$ and $g : X_i' \to X_o'$ be vertical
morphisms (of either value or computation types). Let $R_i : X_i \rel X_i'$ and
$R_o : X_o \rel X_o'$ be horizontal morphisms. We define
%
\[ f \ltsqbisim{R_i}{R_o} g := 
  \exists f', g'.\, f \bisim_{X_i, X_o} f' \ltsq{R_i}{R_o} g' \bisim_{X_i', X_o'} g. \]
%
It remains to verify that the model satisfies the requirements of an extensional
model. 
%
% TODO deriving the original cast rules from the simplified ones in the intensional model
%
The basic idea is that the perturbations in the intensional squares are all
bisimilar to the identity function, and hence they disappear when we perform the
above construction.



\bibliographystyle{plain}
\bibliography{../../paper-new/references}



\end{document}